%% parts/p13/chronology-and-biography.tex --- Part 13
\chapter{Chronology and Biography}\label{sec:p13-chronology-and-biography}

This apparatus is keyed to results and primary works cited elsewhere in the
handbook. Dates identify publication or a documented lecture/result, not an
assertion that a field began in one year. Biographical capsules record the
connection relevant to logic; they are not comprehensive intellectual
biographies.

\section{Selected chronology}
\label{sec:p13-chronology}

\begin{longtable}{@{}p{0.12\textwidth}p{0.80\textwidth}@{}}
\toprule
Date & Milestone \\
\midrule
\endfirsthead
\toprule
Date & Milestone \\
\midrule
\endhead
c.~350 BCE & Aristotle systematizes categorical syllogistic in the
\emph{Prior Analytics}; later traditions distinguish this term logic from
Stoic propositional inference. \\
c.~3rd c. BCE & Stoic logicians, especially Chrysippus, develop analyses of
conditionals and argument schemata closer to propositional logic. \\
1847 & Boole publishes an algebraic calculus of logic \cite{Boole1847}. \\
1879 & Frege's \emph{Begriffsschrift} presents a quantified formal language and
proof system \cite{Frege1879}. \\
1889 & Peano gives a symbolic axiomatization of arithmetic \cite{Peano1889}. \\
1891--97 & Cantor's diagonal arguments and comparison of cardinalities reshape
set theory; the cited 1897 paper supplies a primary record \cite{Cantor1897}. \\
1901--05 & Russell identifies the set-theoretic paradox and publishes the
theory of denoting \cite{Russell1905}. \\
1908 & Zermelo publishes an axiomatization of set theory \cite{Zermelo1908}. \\
1910--13 & Whitehead and Russell publish \emph{Principia Mathematica}
\cite{WhiteheadRussell1910}. \\
1915 & L\"owenheim proves an early downward model-existence theorem for the
relational calculus \cite{Lowenheim1915}. \\
1920s & Skolem sharpens model-theoretic methods and gives a first-order
axiomatization of set theory with the resulting countable-model phenomenon
\cite{Skolem1923}. \\
1928 & Hilbert and Ackermann formulate the decision problem for first-order
logic \cite{HilbertAckermann1928}. \\
1929--30 & Glivenko relates classical and intuitionistic propositional logic;
G\"odel proves first-order completeness \cite{Glivenko1929,Godel1930}. \\
1930 & \L{}ukasiewicz and Tarski publish foundational work on many-valued
propositional calculi \cite{LukasiewiczTarski1930}. \\
1931 & G\"odel proves the first incompleteness theorem
\cite{Godel1931}. \\
1933--35 & Tarski develops a formal conception of truth; Gentzen introduces
natural deduction and sequent calculi \cite{Tarski1983,Gentzen1934a,
Gentzen1935b}. \\
1936--37 & Church and Turing prove undecidability by lambda-definability and
machine computation; Gentzen gives a consistency proof for arithmetic
\cite{Church1936,Turing1937,Gentzen1936}. \\
1944 & Post publishes the structure of recursively enumerable sets and degrees
of unsolvability \cite{Post1944}. \\
1949--50 & Henkin proves completeness for first-order and type-theoretic
calculi \cite{Henkin1949,Henkin1950}. \\
1951--59 & Modal algebra, interpolation, definability, and model theory receive
representation and preservation theorems \cite{JonssonTarski1951,Craig1957,
Beth1953,Tennenbaum1959}. \\
1963--65 & Kripke semantics becomes central for modal and intuitionistic logic;
Robinson introduces resolution \cite{Kripke1963,Kripke1965,Robinson1965}. \\
1969 & Hoare publishes an axiomatic basis for programming; first-order
zero--one laws are proved independently \cite{Hoare1969,GlebskiiEtAl1969}. \\
1971--74 & Cook establishes NP-completeness of SAT; Fagin relates existential
second-order logic to NP \cite{Cook1971,Fagin1974}. \\
1975--77 & Kripke publishes fixed-point truth semantics, Dijkstra weakest
preconditions, Pnueli temporal logic for programs, and Paris--Harrington
independence appears \cite{Kripke1975,Dijkstra1975,Pnueli1977,
ParisHarrington1977}. \\
1981--86 & CTL model checking, symbolic fixed points, BDDs, and descriptive
complexity establish core verification methods \cite{ClarkeEmerson1981,
Vardi1982,Bryant1986,Immerman1986}. \\
1984 & Martin-L\"of's intensional type theory is published in a systematic
form \cite{MartinLof1984}. \\
1987 & Girard introduces linear logic \cite{Girard1987}. \\
1988 & Stable-model semantics gives the foundation of answer-set programming
\cite{GelfondLifschitz1988}. \\
2001--09 & Proof assistants, SAT/SMT technology, separation logic, and verified
systems mature into large proof-producing applications
\cite{ReynoldsSeparation2002,DeMouraBjorner2008,Leroy2009,KleinEtAl2009}. \\
2013--21 & Large collaborative formalizations and the Lean 4 kernel/language
demonstrate new scales of mechanized mathematics
\cite{GonthierEtAl2013,DeMouraUllrich2021}. \\
\bottomrule
\end{longtable}

\section{Biographical capsules}
\label{sec:p13-biographies}

\subsection*{Ancient, early-modern, and algebraic traditions}
The capsules from Aristotle through De Morgan use the historical
reconstructions and source cautions in \cite{KnealeKneale1962}; fragmentary
Stoic evidence is not presented as a surviving modern calculus.

\paragraph{Aristotle (384--322 BCE).}\index{Aristotle}
His syllogistic supplied a systematic theory of valid categorical argument.
Its term-and-figure structure differs from modern quantificational syntax, but
its demand for rule-governed validity is a foundational predecessor.

\paragraph{Chrysippus (c.~279--c.~206 BCE).}\index{Chrysippus}
The Stoic tradition associated with Chrysippus analyzed propositional
connections and argument schemata. Surviving evidence is indirect, so modern
reconstructions must not attribute a complete contemporary calculus verbatim.

\paragraph{Gottfried Wilhelm Leibniz (1646--1716).}\index{Leibniz, Gottfried Wilhelm}
Leibniz pursued a universal characteristic and a calculus of reasoning. The
program anticipated symbolic calculation, although no single surviving system
has the scope of modern predicate logic.

\paragraph{George Boole (1815--1864).}\index{Boole, George}
Boole transformed logical classes and operations into algebraic calculation.
Boolean algebra now bears his name, though its modern equational formulation
was developed through later work as well.

\paragraph{Augustus De Morgan (1806--1871).}\index{De Morgan, Augustus}
De Morgan advanced algebraic logic and the logic of relations. The familiar
duality laws record only one portion of that broader contribution.

\subsection*{Foundations and metamathematics}
The capsules from Frege through Gentzen are keyed to the primary works in
Section~\ref{sec:p13-chronology} and the documented historical synthesis
\cite{GrattanGuinness2000}.

\paragraph{Gottlob Frege (1848--1925).}\index{Frege, Gottlob}
Frege introduced a quantified formal language with explicit binding and proof
rules and used it in a logicist project for arithmetic. His two-dimensional
notation is not the ancestor's syntax reproduced by modern linear notation.

\paragraph{Giuseppe Peano (1858--1932).}\index{Peano, Giuseppe}
Peano developed influential symbolic notation and axiomatizations of
arithmetic. The usual phrase ``Peano arithmetic'' refers to later formal
theories whose exact induction language must still be stated.

\paragraph{Georg Cantor (1845--1918).}\index{Cantor, Georg}
Cantor founded transfinite set theory and introduced diagonal arguments that
separate cardinalities. Diagonalization later became a general method in
logic, computability, and complexity.

\paragraph{Richard Dedekind (1831--1916).}\index{Dedekind, Richard}
Dedekind's structural treatment of the natural numbers and simply infinite
systems helped separate arithmetic structure from a particular realization.

\paragraph{Bertrand Russell (1872--1970).}\index{Russell, Bertrand}
Russell exposed the set-theoretic paradox bearing his name, developed the
theory of descriptions, and coauthored \emph{Principia Mathematica}. These are
distinct achievements, not consequences of one formal device.

\paragraph{Alfred North Whitehead (1861--1947).}\index{Whitehead, Alfred North}
Whitehead coauthored \emph{Principia Mathematica} and contributed to its
ramified type-theoretic program. Later simple type theory is not identical to
that ramified system.

\paragraph{David Hilbert (1862--1943).}\index{Hilbert, David}
Hilbert's program sought formal axiomatization and finitary consistency proofs.
G\"odel's incompleteness theorems delimit particular formal realizations of
that program rather than making metamathematics impossible.

\paragraph{Wilhelm Ackermann (1896--1962).}\index{Ackermann, Wilhelm}
Ackermann worked on the decision problem, proof theory, and recursive
functions. The rapidly growing Ackermann function illustrates the boundary of
primitive recursion.

\paragraph{Leopold L\"owenheim (1878--1957).}\index{Lowenheim@L\"owenheim, Leopold}
L\"owenheim proved an early model-existence theorem for relational logic. The
modern downward L\"owenheim--Skolem theorem combines later formulations and
proof methods.

\paragraph{Thoralf Skolem (1887--1963).}\index{Skolem, Thoralf}
Skolem contributed to model theory, set-theoretic axiomatization, recursive
arithmetic, and Skolem functions. The ``Skolem paradox'' is a contrast between
internal uncountability and external countability, not a contradiction.

\paragraph{Kurt G\"odel (1906--1978).}\index{Godel@G\"odel, Kurt}
G\"odel proved first-order completeness and arithmetical incompleteness and
constructed the constructible universe. Completeness and incompleteness concern
different languages and consequence relations.

\paragraph{Alfred Tarski (1901--1983).}\index{Tarski, Alfred}
Tarski shaped semantic conceptions of truth and consequence, model theory,
relation algebra, and algebraic logic. His truth definitions require an exact
object/metalanguage discipline.

\paragraph{Alonzo Church (1903--1995).}\index{Church, Alonzo}
Church developed the lambda calculus, proved undecidability of first-order
validity, and formulated simple type theory. Church's thesis identifies an
informal effectiveness notion with formal computability and is not itself a
formal theorem.

\paragraph{Alan Turing (1912--1954).}\index{Turing, Alan}
Turing defined machine computation and proved the undecidability of the
halting-style decision problem. His analysis also made the operational content
of an effective procedure unusually explicit.

\paragraph{Gerhard Gentzen (1909--1945).}\index{Gentzen, Gerhard}
Gentzen introduced natural deduction and sequent calculus, proved cut
elimination, and analyzed arithmetic consistency using transfinite induction.
The exact ordinal notation and base theory matter to the final claim.

\subsection*{Mid-century model, proof, and modal theory}
The capsules from Kleene through Craig are keyed to their cited primary
theorems and papers \cite{Kleene1952,Henkin1949,Henkin1950,Post1944,Robinson1965,
Kripke1963,Kripke1965,Kripke1975,Beth1953,Craig1957}.

\paragraph{Stephen Cole Kleene (1909--1994).}\index{Kleene, Stephen Cole}
Kleene developed recursion theory, realizability, the arithmetical hierarchy,
and proof-theoretic methods. His normal-form and recursion theorems anchor
several equivalent presentations of computability.

\paragraph{Leon Henkin (1921--2006).}\index{Henkin, Leon}
Henkin's term-model method proves completeness and extends to type-theoretic
systems under general semantics. Full higher-order semantics has a different
completeness boundary.

\paragraph{Emil Post (1897--1954).}\index{Post, Emil}
Post developed canonical systems, computability models, and the structure of
recursively enumerable degrees. Post's problem initiated the search for
intermediate computably enumerable degrees.

\paragraph{J. Alan Robinson (1923--1974).}\index{Robinson, J. Alan}
Robinson introduced the resolution principle and unification-based automated
deduction for first-order logic. Completeness depends on the declared clause
language, factoring, and equality treatment.

\paragraph{Saul Kripke (1940--2022).}\index{Kripke, Saul}
Kripke developed relational semantics for modal and intuitionistic logics and
fixed-point semantics for truth. These contributions use different formal
structures despite sharing his name.

\paragraph{Evert Beth (1908--1964).}\index{Beth, Evert}
Beth contributed semantic tableaux and the definability theorem: implicit
first-order definability entails explicit definability.

\paragraph{William Craig (1918--2016).}\index{Craig, William}
Craig proved interpolation for first-order logic. Interpolants now serve both
metatheoretic and verification roles, but algorithmic extraction depends on a
proof system.

\subsection*{Nonclassical logic and proof theory}
The capsules from \L{}ukasiewicz through Martin--L\"of use the listed primary
and standard sources \cite{LukasiewiczTarski1930,Brouwer1908,Heyting1930,
Prawitz1965,Girard1972,GirardLafontTaylor1989,Girard1987,MartinLof1984}.

\paragraph{Jan \L{}ukasiewicz (1878--1956).}\index{Lukasiewicz@\L{}ukasiewicz, Jan}
\L{}ukasiewicz developed many-valued calculi and prefix notation and advanced
the historical study of logic. Different \L{}ukasiewicz logics are fixed by
their truth-value algebra.

\paragraph{L. E. J. Brouwer (1881--1966).}\index{Brouwer, L. E. J.}
Brouwer founded intuitionism as a philosophy and practice of mathematics.
Formal intuitionistic logic is one codification and should not be identified
with the whole philosophical program.

\paragraph{Arend Heyting (1898--1980).}\index{Heyting, Arend}
Heyting axiomatized intuitionistic logic and developed its algebraic semantics.
Heyting algebras formalize implication as a residual of meet.

\paragraph{Dag Prawitz (born 1936).}\index{Prawitz, Dag}
Prawitz systematized natural deduction, normalization, and proof-theoretic
semantics. Normalization claims remain calculus-specific.

\paragraph{Jean-Yves Girard (born 1947).}\index{Girard, Jean-Yves}
Girard developed System F, proof-theoretic semantics, and linear logic. Those
systems share proof-theoretic methods but have distinct syntax and metatheory.

\paragraph{Per Martin-L\"of (born 1942).}\index{Martin-Lof@Martin-L\"of, Per}
Martin-L\"of developed constructive type theories in which judgments,
dependent types, and computation rules are foundational. Editions and
intensional/extensional variants must be kept separate.

\subsection*{Logic in computer science and philosophy}
The remaining capsules are keyed respectively to
\cite{Cook1971,Fagin1974,Pnueli1977,Hoare1969,Dijkstra1975,Scott1972Domains,
Quine1937,Quine1948}; birth years are identification data, not claims of current status.

\paragraph{Stephen Cook (born 1939).}\index{Cook, Stephen}
Cook proved the NP-completeness of propositional satisfiability and later
contributed to proof complexity. NP-completeness does not assert that every SAT
instance requires exponential time.

\paragraph{Ronald Fagin (born 1945).}\index{Fagin, Ronald}
Fagin connected existential second-order definability with NP and helped
establish finite model theory and descriptive complexity as interfaces between
logic and computation.

\paragraph{Amir Pnueli (1941--2009).}\index{Pnueli, Amir}
Pnueli introduced temporal logic as a method for reasoning about concurrent
programs. Later LTL and branching-time systems refine distinct path semantics.

\paragraph{C. A. R. Hoare (born 1934).}\index{Hoare, C. A. R.}
Hoare developed axiomatic reasoning about programs, including the triples that
bear his name. Partial and total correctness require different proof
obligations.

\paragraph{Edsger W. Dijkstra (1930--2002).}\index{Dijkstra, Edsger W.}
Dijkstra developed weakest-precondition reasoning and disciplined program
derivation. Weakest preconditions encode termination; weakest liberal
preconditions do not.

\paragraph{Dana Scott (born 1932).}\index{Scott, Dana}
Scott contributed model theory, domain theory, denotational semantics, and
topological interpretations of logic. Scott domains give semantic structure to
recursive definitions and higher-order computation.

\paragraph{W. V. O. Quine (1908--2000).}\index{Quine, W. V. O.}
Quine introduced New Foundations as a set theory and developed the
variable-value criterion of ontological commitment.  The latter applies only
after a theory has been regimented in a declared logic.

\section{Notation and method over time}
\label{sec:p13-notation-history}

\begin{tabularx}{\textwidth}{@{}l X X@{}}
\toprule
Transition & Earlier presentation & Later or concurrent presentation \\
\midrule
categorical inference & subject--predicate terms and figures & quantified formulas with bound variables \\
logical operations & prose and class algebra & truth-functional connectives and algebraic semantics \\
quantification & two-dimensional Fregean scope & linear quantifier notation with parentheses and precedence \\
proof & axioms and informal derivations & natural deduction, sequents, tableaux, resolution, proof terms \\
computation & informal effective procedure & lambda terms, recursive functions, machines, register programs \\
semantics & intended interpretation & structures, valuations, frames, algebras, categories, games \\
mechanization & printed derivation & kernel-checked proof object and versioned dependencies \\
\bottomrule
\end{tabularx}

This table records coexistence, not a simple replacement sequence. Class
algebras, sequent calculi, relational semantics, and proof assistants answer
different questions and remain active when their exact purpose is stated.
